\documentclass[12pt]{article}
\usepackage[T1]{fontenc}
\usepackage[utf8]{inputenc}
\usepackage{lmodern}
\usepackage{textcomp}
\usepackage{enumitem}
\usepackage{geometry}
\geometry{margin=1in}
\begin{document}
\begin{center}
Answer Key - Astronomy - Graduate Level Assessment 3 \
\small (Answers are ordered exactly as in the question paper.)
\end{center}

\section*{Multiple Choice}
\begin{enumerate}[label=\arabic*.]
\item \textbf{Answer:} A
\\ \textit{Options:} A) a dwarf galaxy orbiting the Milky Way.; B) the closest planetary nebula to the Earth.; C) a bright star cluster discovered by Magellan.; D) the outer arm of the Milky Way named after Magellan.
\\ \textit{Note:} The Large Magellanic Cloud is classified as a dwarf galaxy because it is a smaller, irregular galaxy that orbits the Milky Way, demonstrating the gravitational interactions typical of satellite galaxies in our local group.
\item \textbf{Answer:} A
\\ \textit{Options:} A) size of the planet; B) presence of an atmosphere; C) distance from the sun; D) rotation period
\\ \textit{Note:} The size of a planet is the most important factor in determining its volcanic and tectonic activity because larger planets retain heat longer, which enables geological processes to continue, whereas smaller planets cool more quickly, leading to reduced or halted volcanic and tectonic activity.
\item \textbf{Answer:} A
\\ \textit{Options:} A) Its rotation is too slow.; B) It has too thick an atmosphere.; C) It is too large.; D) It does not have a metallic core.
\\ \textit{Note:} Venus has a slow rotation period of about 243 Earth days, which hinders the dynamo effect necessary for generating a strong magnetic field, as this process requires a fast rotation to facilitate the movement of molten iron in its core.
\item \textbf{Answer:} B
\\ \textit{Options:} A) 5; B) 7; C) 9; D) 12
\\ \textit{Note:} The Pleiades is traditionally noted for its visibility of seven prominent stars to the naked eye, which has led to its significance in various mythologies and cultural narratives throughout history.
\item \textbf{Answer:} C
\\ \textit{Options:} A) Because the molecules that compose the Earth's atmosphere have a blue-ish color.; B) Because the sky reflects the color of the Earth's oceans.; C) Because the atmosphere preferentially scatters short wavelengths.; D) Because the Earth's atmosphere preferentially absorbs all other colors.
\\ \textit{Note:} The sky appears blue because Rayleigh scattering causes shorter wavelengths of light, such as blue, to be scattered more effectively by the gases and particles in the Earth's atmosphere than longer wavelengths like red.
\end{enumerate}

\section*{True / False}
\begin{enumerate}[label=\arabic*.]
\setcounter{enumi}{5}
\item \textbf{Answer:} False
\\ \textit{Options:} A) True; B) False
\\ \textit{Note:} The statement is false because hydrogen is the most abundant element in the atmospheres of the Jovian planets, significantly surpassing helium in terms of abundance.
\item \textbf{Answer:} False
\\ \textit{Options:} A) True; B) False
\\ \textit{Note:} The statement is false because, while Mars is indeed farther from the Sun than Earth, its more extreme seasonal variations are primarily due to its axial tilt and thin atmosphere rather than its distance from the Sun, which results in cooler average temperatures but not necessarily more extreme seasonal changes.
\item \textbf{Answer:} True
\\ \textit{Options:} A) True; B) False
\\ \textit{Note:} The atmosphere on Mars is composed primarily of carbon dioxide but is only about 1\% as dense as Earth's atmosphere, which significantly reduces its capacity to retain heat and thus limits the effectiveness of the greenhouse effect.
\item \textbf{Answer:} True
\\ \textit{Options:} A) True; B) False
\\ \textit{Note:} Moons exert gravitational forces that help maintain the structure and stability of ring systems by influencing particle orbits, creating gravitational interactions at ring edges, and facilitating the formation of gaps through their gravitational perturbations.
\item \textbf{Answer:} False
\\ \textit{Options:} A) True; B) False
\\ \textit{Note:} The statement is false because if the angular separation of two stars is small, they may be too close together to be resolved as distinct objects by the human eye, leading to the perception of a single point of light instead.
\end{enumerate}

\section*{Long Form Response}
\begin{enumerate}[label=\arabic*.]
\setcounter{enumi}{10}
\item \textbf{Answer:} Orbital resonance occurs when two orbiting bodies exert regular, periodic gravitational influence on each other, often leading to stable orbital configurations and influencing their motion over time. This phenomenon is crucial in understanding the dynamics of celestial bodies, such as the resonances found in the orbits of moons around giant planets. In the case of the small Jovian moons, their breaking to form ring materials does not primarily result from orbital resonance but rather from tidal forces and collisions. These moons are often subjected to intense gravitational interactions and tidal stresses that exceed their structural integrity, leading to fragmentation. Thus, while orbital resonance plays a significant role in the stability of larger celestial bodies, the disintegration of small moons is more directly linked to their susceptibility to these disruptive forces rather than resonant interactions.
\\ \textit{Note:} correct because it accurately identifies that orbital resonance involves periodic gravitational interactions that stabilize orbits, while the destruction of small Jovian moons is primarily caused by tidal forces rather than these resonant interactions.
\item \textbf{Answer:} The significance of black bodies in astronomy lies in their role as idealized physical objects that absorb all electromagnetic radiation, thus serving as a benchmark for understanding stellar radiation. A perfect black body not only absorbs but also emits radiation in a characteristic spectrum solely dependent on its temperature, described by Planck's law. This concept aids astronomers in interpreting the spectral emissions of stars and other celestial bodies, allowing for the determination of their temperatures, compositions, and distances. By comparing observed spectra to the theoretical models of black body radiation, researchers can derive crucial astrophysical parameters, enhancing our comprehension of cosmic phenomena and the underlying physics governing stellar and galactic evolution.
\\ \textit{Note:} correct because it accurately defines a black body as an idealized object that absorbs and emits radiation according to its temperature, highlighting its significance in astronomy for interpreting stellar radiation through Planck's law.
\end{enumerate}

\vfill
\noindent\textit{End of Answer Key}
\end{document}